\chapter{交互、编队、战报与手册}
\section{区域与返回上下文}
主导航为基地、编队、战役、世界、指挥官、战报。工厂和科研通过基地对应建筑进入。资源区选资源场与仓库；工业区分别进两制造厂、改装与维修；指挥科研区连接HQ、科技树和学院。车间使用同一绘制函数，但选择的factory标识决定报价与队列，不能混成共享生产线。

返回按钮统一靠页面右上或区域既定位置，导航栈还原选中建筑和车间模式。进入新档清导航、战前上下文和原地清单滚动；日常tick则不重置操作位置。部分全屏工具与档案仍分页，主基地右栏建设/科研是原生滚动列表。

\section{编队、预设与连续部署}
六格计划按同车型共享库存，不能每格各自把总库存算一次。加载预设按阵位顺序分配：每格先钳制统率，再钳制该车剩余库存，记录stockShortage和capShortage。原命名预设保留计划数，当前出战取可用最大，库存不足提示各格缺额。

候选车型部署后仍选中，左侧切换阵位即可再次部署；每次重新算剩余库存。点击最大战力按真实库存拆成最多六组选择，最高兵阶模式优先阶级，再比总分；前排优先放更高总基础生命的已选组，不改变战力总和。

战前部署允许草稿调整、查看敌军属性、取消返回；正式确认同时保存该次formation并出战。连续下一关沿用当前可用编队直接battle/dungeon，不再强制进备战；空编队阻止离开原结算，待响应重复点击只提交一次。
\src{src/core/engine.ts}{formationAvailability；battleFormation}
\src{src/core/power.ts}{arrangedFormation}

\section{战报与归队账本}
400条预算由战斗和归队收据共享，按时间倒序、序号与ID稳定排序。裁掉战报不清除首通与领奖事实。分类含主线、核心、攻矿、突袭、演习、采集归队和突袭归队，另按成功/失败过滤。

世界攻矿战报回答“是否打赢与损失多少”；归队收据回答“何时回来、实际入库多少、超仓丢弃多少”。二者通过marchId/battleId关联，不能因为找不到当前march就推断旧货物已入库；旧收据缺stored时明确是历史货物记录而非可证实的入库账。

归队success不仅看outcome=returned，还看该次battleWon；败战幸存部队回城仍是失败任务。safe矿没有战斗也可以成功收据。战报页进入时复制列表，用户手动刷新以免tick插入新记录造成点击位置移动。

\section{休整结算浮窗}
直接点击休整推进时间，返回资源净变化、自产/仓容限制、建筑/科研完成、制造/改装/维修交付、归队及新增战报。浮窗保留资源与奖励图标，点击任意处关闭，没有再次确认和操作按钮。

资源净变化包含自产、归队和结算期间其他钱包变化，不等于理论速率乘时间。produced用总produce计数减converted差额区分新造与改装；战损来自不可变行军历史，不能误当当前待修库存。
\src{src/core/overview.ts}{progressSummary；transportStatus}
\src{native/scripts/rest_summary.gd}{show_summary}

\section{图书馆与显示规则}
图书馆当前40篇，按分类和关键字筛选，每篇有稳定ID与内容，目录由桥接下发。阅读器按文章ID记位置，返回战斗时恢复暂停状态；相关入口可跳科技或说明页面。版本标识从package版本生成，避免升级后手册仍显示旧版本。

资源格式使用十进制K/M/G/T，每档1000，最多两位并截断而非四舍五入，避免显示比实际更多的可用材料。消费不足可标红，奖励使用非消费语义，不以当前余额判断奖励颜色。普通命中、暴击、闪避与未破防采用不同字号/颜色与标签。
\src{src/core/archive.ts}{trimArchive；archiveRows}
\src{src/core/library.ts}{fieldLibrary；searchLibrary}
\src{src/core/format.ts}{compactAmount}
